% !TEX root = ../main.tex

\newcommand{\subseteqfin}{\subseteq_\text{\tiny fin}}
\newcommand{\setto}[1]{\{1\ldots #1\}}
\newcommand{\updinstr}[1]{\textbf{update}(#1)}
\newcommand{\ztinstr}[1]{\textbf{zero-test}(#1)}
\newcommand{\mach}{\mathcal{M}}
\newcommand{\Instr}{\mathcal{I}}
\newcommand{\powset}[1]{\mathcal{P}(#1)}
\newcommand{\wrt}{w.r.t.\xspace}
\newcommand{\ie}{i.e.\xspace}
\newcommand{\eg}{e.g.\xspace}
\newcommand{\pn}{\text{\sc pn}\xspace}
\newcommand{\vas}{\text{\sc vas}\xspace}
\newcommand{\vass}{\text{\sc vass}\xspace}
\newcommand{\cm}{\text{\sc cm}\xspace}

\newcommand{\dvass}{data \text{\sc vass}\xspace}
\newcommand{\onevas}{1-\text{\sc vas}\xspace}
\newcommand{\onevass}{1-\text{\sc vass}\xspace}
\newcommand{\donevas}{1-\text{\sc dvas}\xspace}
\newcommand{\donevass}{1-\text{\sc dvass}\xspace}
\newcommand{\donevasr}{1-\text{\sc dvasr}\xspace}
\newcommand{\donepn}{1-\text{\sc dpn}\xspace}

\newcommand{\cycg}[1]{\textsc{Z}_{#1}}
\newcommand{\symg}[1]{\textsc{S}_{#1}}
\newcommand{\sym}{\text{Sym}\xspace}

\newcommand{\dtwovas}{2-\text{\sc dvas}\xspace}
\newcommand{\dtwovass}{2-\text{\sc dvass}\xspace}
\newcommand{\dtwopn}{2-\text{\sc dpn}\xspace}


\newcommand{\vr}[1]{\mathbf{#1}}%{\vect{#1}}
\newcommand{\setof}[2]{\{\,#1 \ \mid \ #2\,\}}
\newcommand{\aut}[1]{\textsc{Perm}(#1)}
\newcommand{\step}{\xlongrightarrow[]{\quad}}
%\newcommand{\step}{\longrightarrow}
\newcommand\xstep[1]{\xrightarrow{#1}}
%\newcommand{\xstep}[1]{\stackrel{#1}{\step}}
\newcommand{\Mywlog}{W.l.o.g.\xspace}
\newcommand{\mywlog}{w.l.o.g.\xspace}
\newcommand{\norm}[1]%{\textsc{norm}(#1)}
{\lVert#1\rVert}
\newcommand{\width}[1]{\textsc{width}(#1)}
\newcommand{\height}[1]{\textsc{height}(#1)}%{\lVert#1\rVert}
\newcommand{\findim}[2]{{#1}_{#2}}
\newcommand{\decomp}{decomposable\xspace}
\newcommand{\indecomp}{indecomposable\xspace}
\newcommand{\Decomp}{Decomposable\xspace}
\newcommand{\Indecomp}{Indecomposable\xspace}
\newcommand{\trace}{trace\xspace}
\newcommand{\prefsum}[2]{#1[#2]}
\newcommand{\reachprob}[1]{#1 \textsc{reachability}\xspace}
\newcommand\lin[2]{\mathcal{L}\left(#1;#2\right)} %linear set


\newcommand{\np}{{\sc NP}\xspace}
\newcommand{\nl}{{\sc NL}\xspace}
\newcommand{\pspace}{{\sc PSpace}\xspace}
\newcommand{\expspace}{{\sc ExpSpace}\xspace}
\newcommand{\Ackermann}{{{\sc Ackermann}}\xspace}
\newcommand{\para}[1]{%\vspace{-1mm}
\subparagraph*{#1.}}
\newcommand{\tuple}[1]{\ensuremath{\langle #1 \rangle}\xspace}

%%% Todo notes
\newcommand{\sltodo}{\textcolor{blue}{TODO}}
\newcommand{\sla}[1]{\textcolor{blue}{SL: #1}}
\newcommand{\mei}[1]{\textcolor{magenta}{MK: #1}}
\newcommand{\vrm}[1]{\textcolor{teal}{VM: #1}\par}
\newcommand{\ml}[1]{\todo[color=cyan]{ML: {#1}}}
\newcommand{\rp}[1]{\todo[color=green]{RP: {#1}}}
\newcommand{\mlinline}[1]{\todo[color=cyan,inline]{ML: {#1}}}
\newcommand{\red}[1]{\textcolor{red}{#1}}

%%% Sets
\newcommand{\R}{\mathbb{R}}
\newcommand{\Z}{\mathbb{Z}}
\newcommand{\N}{\mathbb{N}}
\newcommand{\Npos}{\mathbb{N}_{>0}}
\newcommand{\Q}{\mathbb{Q}}
\newcommand{\Time}{\mathbb{T}}
\newcommand{\Qpos}{\Q_{\geq 0}}
\newcommand{\Rpos}{\R_{\geq 0}}
\newcommand{\atoms}{\mathbb{A}}

\newcommand{\finto}{\to_{\text{fin}}}
\newcommand{\finsubseteq}{\subseteq_{\text{fin}}}
\newcommand{\V}{\mathcal{V}}

\newcommand{\A}{\mathcal{A}}
\newcommand{\B}{\mathcal{B}}
\newcommand{\C}{\mathcal{C}}
\newcommand{\mysize}[1]{|#1|}

%%% TIKZ
\usetikzlibrary{automata,positioning,shadows}
\usetikzlibrary{calc, shapes, backgrounds,arrows}
\usetikzlibrary{chains, decorations.pathmorphing}
\usetikzlibrary{positioning,decorations.pathreplacing, automata}
\usetikzlibrary{shapes.misc}

\tikzset{every loop/.style={looseness=7}, >=latex}
\tikzset{every picture/.style={>=latex}}
\tikzstyle{PlayerMin}=[draw,circle,minimum size=4mm,inner sep=1.5pt]
\tikzstyle{PlayerMax}=[draw,rectangle,minimum size=7mm,inner sep=1.5pt]
\tikzstyle{Player}=[draw,diamond,minimum size=7mm,inner sep=1.5pt]
\tikzstyle{target}=[circle, minimum size=1mm,inner sep=-2pt]
\tikzstyle{PlayerMinmin}=[draw,circle, minimum size=1.5mm,inner sep=0pt]
\tikzstyle{PlayerMaxmin}=[draw,rectangle,minimum size=1.5mm,inner sep=0pt]
\tikzstyle{PlayerMinsmall}=[draw,circle, minimum size=3mm,inner sep=0pt]
\tikzstyle{PlayerMaxsmall}=[draw,rectangle,minimum size=3mm,inner sep=0pt]
\tikzstyle{leaf}=[draw,diamond,minimum size=7mm,inner sep=1.5pt]

\tikzstyle{strat} =[minimum width=0.1cm,line width=0.01mm,draw=none]
\tikzstyle{vecArrow} = [decoration={markings,mark=at position
1 with {\arrow[scale=.5,>=latex]{>}}}%,postaction={decorate}
]



\theoremstyle{definition}
\newtheorem{fact}{Fact}
\newtheorem{question}{Question}

